360
6 Integration
\medskip
\begin{enumerate}
    \setcounter{enumi}{3}
    \item a) Show that if $f \in \mathcal{R}[a, b]$, then $| f |^p \in \mathcal{R}[a, b]$ for $p \ge 0$.
    \medskip
    b) Starting from Hölder's inequality for sums, obtain Hölder's inequality for
    integrals:$^6$
    \[
    \left| \int_a^b (f \cdot g)(x)\,dx \right| \le \left( \int_a^b |f|^p(x)\,dx \right)^{1/p} \left( \int_a^b |g|^q(x)\,dx \right)^{1/q},
    \]
    if $f, g \in \mathcal{R}[a, b]$, $p > 1$, $q > 1$, and $\frac{1}{p} + \frac{1}{q} = 1$.
    \medskip
    c) Starting from Minkowski's inequality for sums, obtain Minkowski's inequal-
    ity for integrals:
    \[
    \left( \int_a^b |f + g|^p(x)\,dx \right)^{1/p} \le \left( \int_a^b |f|^p(x)\,dx \right)^{1/p} + \left( \int_a^b |g|^p(x)\,dx \right)^{1/p},
    \]
    if $f, g \in \mathcal{R}[a, b]$ and $p \ge 1$. Show that this inequality reverses direction if $0 <
    p < 1$.
    \medskip
    d) Verify that if $f$ is a continuous convex function on $\mathbb{R}$ and $\varphi$ an arbitrary
    continuous function on $\mathbb{R}$, then Jensen's inequality
    \[
    f\left(\frac{1}{c}\int_0^c \varphi(t)\,dt\right) \le \frac{1}{c}\int_0^c f(\varphi(t))\,dt
    \]
    holds for $c\neq 0$.
\end{enumerate}
\medskip
\noindent \textbf{6.3 The Integral and the Derivative}
\medskip
\noindent \textbf{6.3.1 The Integral and the Primitive}
\medskip
Let $f$ be a Riemann-integrable function on a closed interval $[a, b]$. On this interval
let us consider the function
\[
F(x) = \int_a^x f(t)\,dt,
\tag{6.40}
\]
often called an integral with variable upper limit.
\medskip
Since $f \in \mathcal{R}[a, b]$, it follows that $f|_{[a,x]} \in \mathcal{R}[a, x]$ if $[a,x] \subset [a, b]$; therefore
the function $x \to F(x)$ is unambiguously defined for $x \in [a, b]$.
\medskip
\noindent $^6$The algebraic Hölder inequality for $p = q = 2$ was first obtained in 1821 by Cauchy and bears his
name. Hölder's inequality for integrals with $p = q = 2$ was first discovered in 1859 by the Russian
mathematician B.Ya. Bunyakovskii (1804–1889). This important integral inequality (in the case
$p = q = 2$) is called Bunyakovskii's inequality or the Cauchy–Bunyakovskii inequality. One also
some-times sees the less accurate name “Schwarz inequality" after the German mathematician
H.K.A. Schwarz (1843–1921), in whose work it appeared in 1884.